\chapter{Quantum Circuit Transpilation and Hardware Resource Profiles}\label{app:transpilation_profiles}

\notatfm{Empirical gate metrics extracted from direct compilation passes in Qiskit 1.x / AerSimulator under seed 100.}
This appendix provides the comprehensive transpilation profiles, physical gate breakdowns, and asymptotic hardware resource scalings for the quantum circuits developed throughout this thesis. All metrics presented herein were obtained empirically by compiling the quantum circuits within the Qiskit framework targeting realistic superconducting quantum processor gate bases.

\section{Transpilation Toolchain and Basis Set}\label{sec:app_transpiler_env}

The compilation pipeline translates high-level parameterized quantum Bayesian circuits and amplitude amplification operators into concrete gate instructions executable on physical quantum processing units (QPUs) or full-state density matrix simulators. 

\subsection{Execution Environment and Compiler Directives}
The circuit synthesis, transpilation passes, and resource accounting were executed under the following standard configuration:
\begin{itemize}[leftmargin=*]
    \item \textbf{Quantum SDK:} Qiskit framework ($\ge 1.0$) with high-performance C++ backend simulation via \texttt{qiskit-aer} (\texttt{AerSimulator}).
    \item \textbf{Compiler Optimization:} Pass Manager optimization level set to \texttt{optimization\_level=1}, applying standard single-qubit gate cancellations, commutativity-based gate mergers, and topological layout mapping.
    \item \textbf{Transpiler Seed:} Deterministic synthesis ensured via \texttt{seed\_transpiler=100}.
    \item \textbf{Target Basis Set:} Universal gate set including single-qubit rotations and entangling operations:
    \begin{equation}
    \mathcal{B}_{\mathrm{target}} = \left\{ \mathrm{CX}, \mathrm{CCX}, \mathrm{U}_1, \mathrm{U}_2, \mathrm{U}_3, \mathrm{RY}, \mathrm{CRZ}, \mathrm{CU}, \mathrm{CP}, \mathrm{H}, \mathrm{X}, \mathrm{Z}, \mathrm{MCPHASE}, \mathrm{MCX}, \mathrm{SWAP} \right\}.
    \end{equation}
    \item \textbf{Subcircuit Modularity:} All conditional probability subcircuits ($\mathcal{A}$, $S_\chi$, $\mathcal{Q}$, and $\mathrm{IQFT}^\dagger$) follow strict barrier-free specifications, ensuring that converter calls (\texttt{to\_gate()}) preserve unhindered peephole optimizations and recursive unrolling during transpilation.
\end{itemize}

\section{Master 17-Qubit QAE Circuit Profile}\label{sec:app_master_qae_profile}

The master Quantum Amplitude Estimation (QAE) architecture \cite{brassard2002} formalizes the full K2-18b astrobiological inference task over $17$ qubits, structured into three functionally dedicated quantum registers:
\begin{enumerate}[leftmargin=*]
    \item \textbf{Evaluation Register ($\mathbf{q}_E$, 5 qubits, $q_0 \dots q_4$):} Discretizes the phase readout register for the $m=5$ Quantum Phase Estimation stage, providing an angular resolution of $\Delta\theta = \pi / 2^5 \approx 0.098175\,\mathrm{rad}$.
    \item \textbf{Astrobiological State Register ($\mathbf{q}_S$, 9 qubits, $q_5 \dots q_{13}$):} Directly encodes the nine physical and biosignature nodes $\mathbf{X} = \{X_0, \dots, X_8\}$ of the K2-18b Bayesian network via amplitude encoding.
    \item \textbf{Ancilla Workspace Register ($\mathbf{q}_A$, 3 qubits, $q_{14} \dots q_{16}$):} Facilitates clean uncomputation and multi-controlled gate decompositions for the phase oracle $S_\chi$ and zero-state reflection diffuser $S_0$.
    \item \textbf{Classical Readout Register ($\mathbf{c}_R$, 5 classical bits):} Stores the projective measurement outcomes of $\mathbf{q}_E$ upon execution of the Inverse Quantum Fourier Transform ($\mathrm{IQFT}^\dagger$).
\end{enumerate}

\subsection{Subcircuit Component Metrics}
Prior to global circuit assembly, each fundamental operator subcircuit was transpiled individually under \texttt{optimization\_level=1} to establish baseline compilation footprints:
\begin{itemize}[leftmargin=*]
    \item \textbf{State Preparation Operator $\mathcal{A}$ ($12$ operational qubits: $9$ state + $3$ ancilla):}
    \begin{itemize}
        \item Transpiled Depth: $33$
        \item Total Gate Count (Size): $47$
        \item Gate Breakdown: $14\times\mathrm{X}$, $12\times\mathrm{CU}$, $11\times\mathrm{CRY}$, $8\times\mathrm{CX}$, $2\times\mathrm{U}_3$.
    \end{itemize}
    \item \textbf{Phase Oracle $S_\chi$ ($12$ operational qubits):}
    \begin{itemize}
        \item Transpiled Depth: $1$
        \item Total Gate Count (Size): $1$
        \item Gate Breakdown: $1\times\mathrm{Z}$ (marking target state on qubit $q_8$).
    \end{itemize}
    \item \textbf{Grover Iterate Operator $\mathcal{Q} = -\mathcal{A} S_0 \mathcal{A}^\dagger S_\chi$ ($12$ operational qubits):}
    \begin{itemize}
        \item Transpiled Depth: $67$
        \item Total Gate Count (Size): $112$
        \item Gate Breakdown: $42\times\mathrm{X}$, $24\times\mathrm{CU}$, $22\times\mathrm{CRY}$, $16\times\mathrm{CX}$, $4\times\mathrm{U}_3$, $2\times\mathrm{U}_2$, $1\times\mathrm{Z}$, $1\times\mathrm{MCX}$.
    \end{itemize}
    \item \textbf{5-Qubit Inverse Quantum Fourier Transform ($\mathrm{IQFT}^\dagger$, $5$ evaluation qubits):}
    \begin{itemize}
        \item Transpiled Depth: $10$
        \item Total Gate Count (Size): $17$
        \item Gate Breakdown: $10\times\mathrm{CP}$, $5\times\mathrm{H}$, $2\times\mathrm{SWAP}$.
    \end{itemize}
\end{itemize}

\subsection{Full Master Circuit Assembly and Gate Breakdown}
The complete master QAE circuit cascades $\sum_{j=0}^{4} 2^j = 31$ controlled-Grover iterates $C\text{-}\mathcal{Q}^{2^j}$ across the evaluation register, bounded by a layer of Hadamard gates and finalized by $\mathrm{IQFT}^\dagger$ and classical measurement.

Transpilation of the full $17$-qubit circuit yields an unrolled circuit depth of $\mathbf{52{,}393}$ and a total gate volume of $\mathbf{53{,}456}$ operations. Table \ref{tab:master_qae_gates} details the complete, exact gate inventory extracted from the compiler pass.

\begin{table}[htbp]
\centering
\caption{\textbf{Transpiled Gate Breakdown for the Master 17-Qubit QAE Circuit.}}
\label{tab:master_qae_gates}
\small
\begin{tabular}{llrr}
\toprule
\textbf{Quantum Gate Type} & \textbf{Qiskit Instruction} & \textbf{Absolute Count} & \textbf{Relative Share (\%)} \\
\midrule
Multi-Controlled Phase & \texttt{mcphase} & $17{,}515$ & $32.765\%$ \\
Toffoli (Controlled-Controlled-NOT) & \texttt{ccx} & $17{,}112$ & $32.011\%$ \\
Single-Qubit General Rotation & \texttt{u2} & $8{,}680$ & $16.238\%$ \\
Controlled-NOT (CNOT) & \texttt{cx} & $5{,}991$ & $11.207\%$ \\
Controlled-Unitary & \texttt{cu} & $2{,}988$ & $5.590\%$ \\
Controlled-$R_Z$ Rotation & \texttt{crz} & $1{,}054$ & $1.972\%$ \\
Hadamard & \texttt{h} & $72$ & $0.135\%$ \\
Pauli-$X$ & \texttt{x} & $14$ & $0.026\%$ \\
Controlled-$R_Y$ Rotation & \texttt{cry} & $11$ & $0.021\%$ \\
Controlled-Phase & \texttt{cp} & $10$ & $0.019\%$ \\
Projective Measurement & \texttt{measure} & $5$ & $0.009\%$ \\
Single-Qubit General Rotation & \texttt{u3} & $2$ & $0.004\%$ \\
Qubit Swap & \texttt{swap} & $2$ & $0.004\%$ \\
\midrule
\textbf{Total Operations} & & $\mathbf{53{,}456}$ & $\mathbf{100.000\%}$ \\
\bottomrule
\end{tabular}
\end{table}

\notatfm{Multi-qubit controlled gates (\texttt{mcphase}, \texttt{ccx}, \texttt{cx}) constitute $>75\%$ of the master circuit volume.}
As shown in Table \ref{tab:master_qae_gates}, multi-controlled phase gates (\texttt{mcphase}, $32.77\%$) and Toffoli gates (\texttt{ccx}, $32.01\%$) dominate the compiled instruction stream. This concentration arises directly from the sequential conditioning of the 9-node DAG and the multi-qubit reflection operator $S_0 = 2\ket{0}\bra{0} - I$, which requires high-order entanglement across the state and ancilla registers.

\subsection{Native Two-Qubit CNOT Decomposition Overhead}\label{subsec:app_cnot_decomposition_overhead}
It is important to distinguish between the abstract gate counts reported in Table \ref{tab:master_qae_gates} (which retain multi-controlled unitary instructions in the universal basis $\mathcal{B}_{\mathrm{target}}$) and the fully decomposed physical CNOT depth. Standard unitary synthesis of a 3-qubit Toffoli gate (\texttt{ccx}) requires at least 6 CNOT gates and 7 single-qubit rotations \cite{barenco1995}. Consequently, unrolling the $17{,}112$ Toffoli operations alone contributes $17{,}112 \times 6 = 102{,}672$ CNOT operations to the physical execution graph, while decomposing the $17{,}515$ multi-controlled phase gates (\texttt{mcphase}) introduces further entangling layers. 

This brings the total native two-qubit CNOT count well beyond $105{,}000$ operations when compiled without auxiliary borrowed-ancilla networks. This substantial compilation footprint provides the empirical justification for classifying canonical QAE as an exclusive Fault-Tolerant Quantum Computing (FTQC) blueprint, directly motivating the transition to Iterative QAE (IQAE) for near-term physical QPU viability.

\section{NISQ Kernel Transpilation and ZNE Folding Progression}\label{sec:app_nisq_kernel_profile}

To benchmark performance on noisy intermediate-scale hardware without fault-tolerant logical ancillae, the 5-qubit astrobiological inference kernel was subjected to Zero-Noise Extrapolation (ZNE) via unitary gate folding \cite{temme2017, li2017, endo2018}:
\begin{equation}
\mathcal{U} \;\longrightarrow\; \mathcal{U} \left(\mathcal{U}^\dagger \mathcal{U}\right)^k, \quad \text{with noise scaling factor } \lambda = 2k + 1 \in \{1, 3, 5\}.
\end{equation}

\subsection{Physical Noise Environment}
The kernel was transpiled and simulated against a realistic superconducting transmon noise model calibrated to the following parameters:
\begin{itemize}[leftmargin=*]
    \item Coherence lifetimes: Longitudinal relaxation $T_1 = 150\,\mu\mathrm{s}$, transverse dephasing $T_2 = 120\,\mu\mathrm{s}$.
    \item Gate durations: Single-qubit pulse duration $t_{1q} = 35\,\mathrm{ns}$, two-qubit cross-resonance duration $t_{2q} = 300\,\mathrm{ns}$.
    \item Gate Pauli errors: Single-qubit depolarizing error $p_{1q} = 0.08\%$, two-qubit depolarizing error $p_{2q} = 1.20\%$.
    \item Measurement assignment errors: Readout probabilities $P(1\mid 0) = 1.50\%$, $P(0\mid 1) = 1.00\%$.
\end{itemize}

\subsection{Empirical ZNE Resource Progression}
Table \ref{tab:zne_folding_scaling} presents the exact uncompiled and transpiled resource profiles across the unitary folding ladder $\lambda \in \{1, 3, 5\}$.

\begin{table}[htbp]
\centering
\caption{\textbf{Hardware Compilation Profile across ZNE Unitary Folding Factors ($\lambda \in \{1, 3, 5\}$).}}
\label{tab:zne_folding_scaling}
\small
\begin{tabular}{lrrr}
\toprule
\textbf{Metric / Operation} & \textbf{Baseline ($\lambda = 1$)} & \textbf{Folded ($\lambda = 3$)} & \textbf{Folded ($\lambda = 5$)} \\
\midrule
\multicolumn{4}{l}{\textit{Abstract Circuit Geometry}} \\
Uncompiled Circuit Depth & $10$ & $28$ & $46$ \\
Uncompiled Total Gate Count & $21$ & $61$ & $101$ \\
\midrule
\multicolumn{4}{l}{\textit{Transpiled Physical Profile (\texttt{optimization\_level=1}, \texttt{seed=100})}} \\
Transpiled Circuit Depth & $\mathbf{39}$ & $\mathbf{118}$ & $\mathbf{197}$ \\
Transpiled Total Gate Count & $\mathbf{57}$ & $\mathbf{156}$ & $\mathbf{255}$ \\
\midrule
\multicolumn{4}{l}{\textit{Detailed Transpiled Gate Inventory}} \\
Controlled-NOT (\texttt{cx}) & $23$ & $69$ & $115$ \\
Single-Qubit Rotation (\texttt{u3}) & $22$ & $64$ & $106$ \\
Pauli-$Y$ Rotation (\texttt{ry}) & $4$ & $8$ & $12$ \\
Pauli-$X$ (\texttt{x}) & $3$ & $6$ & $9$ \\
Single-Qubit Rotation (\texttt{u2}) & $3$ & $4$ & $5$ \\
Hadamard (\texttt{h}) & $1$ & $2$ & $3$ \\
Phase Rotation (\texttt{u1}) & $0$ & $2$ & $4$ \\
Projective Readout (\texttt{measure}) & $1$ & $1$ & $1$ \\
\midrule
\textbf{Total Transpiled Operations} & $\mathbf{57}$ & $\mathbf{156}$ & $\mathbf{255}$ \\
\bottomrule
\end{tabular}
\end{table}

\subsection{Linearity of Two-Qubit Error Amplification}
\notatfm{CNOT counts scale strictly as $23 \cdot \lambda$, guaranteeing linear decoherence amplification for Richardson extrapolation.}
A critical requirement for the validity of polynomial Zero-Noise Extrapolation is that gate folding amplifies physical noise strictly in proportion to $\lambda$ without introducing non-linear structural artifacts during transpilation. 

The empirical metrics in Table \ref{tab:zne_folding_scaling} confirm this linear scaling:
\begin{itemize}[leftmargin=*]
    \item \textbf{Entangling CNOT Gate Count:} Compiles to exactly $N_{\mathrm{CX}}(\lambda) = 23 \times \lambda$, yielding $23$, $69$, and $115$ CNOTs for $\lambda = 1, 3, 5$ respectively.
    \item \textbf{Circuit Depth Progression:} Advances from $39 \to 118 \to 197$, exhibiting a steady incremental depth of $\Delta D \approx 79$ per unitary inversion cycle $(\mathcal{U}^\dagger \mathcal{U})$.
    \item \textbf{Total Gate Volume:} Advances from $57 \to 156 \to 255$, corresponding to an exact increment of $\Delta N_{\mathrm{ops}} = 99$ operations per fold cycle ($46$ CNOTs plus $53$ single-qubit rotations).
\end{itemize}
This strict algebraic proportionality ensures that the empirical noise parameter matches the theoretical assumption underlying Richardson polynomial extrapolation, enabling the $86.44\%$ error mitigation efficiency demonstrated in Chapter \ref{chap:system_architecture}.

\section{Asymptotic Hardware Scaling Across Evaluation Widths}\label{sec:app_scaling_analysis}

In canonical QAE, the precision of the estimated phase $\tilde{\theta}$ is governed by the evaluation register width $n_E$. The phase estimation grid discretizes the interval $[0, \pi]$ into intervals of width:
\begin{equation}
\Delta\theta = \frac{\pi}{2^{n_E}},
\end{equation}
which dictates the minimum resolvable separation between rare-event probability amplitudes. However, the total number of controlled-Grover iterations scales exponentially with $n_E$:
\begin{equation}
N_Q = \sum_{j=0}^{n_E - 1} 2^j = 2^{n_E} - 1.
\end{equation}

Table \ref{tab:qae_scaling_registers} compiles the empirical and projected hardware resource requirements as the evaluation precision increases from $n_E = 3$ to $n_E = 10$. The projected gate volumes follow the analytical scaling model $\text{Estimated Gates} \approx N_Q \times 1{,}724 = (2^{n_E} - 1) \times 1{,}724$, where each controlled Grover iterate accounts for approximately $1{,}724$ compiled physical instructions. For the baseline $n_E = 5$, this closed-form model yields $53{,}444$ operations, matching the empirically transpiled master circuit total of $53{,}456$ within $0.02\%$ (with the residual difference attributable to the $\mathrm{IQFT}^\dagger$ permutation network and classical readout measurement operations).

\begin{table}[htbp]
\centering
\caption{\textbf{Asymptotic Hardware Scaling across Evaluation Register Precision ($n_E \in \{3, \dots, 10\}$).}}
\label{tab:qae_scaling_registers}
\small
\begin{tabular}{rrrrr}
\toprule
\textbf{Register Size ($n_E$)} & \textbf{Total Iterates ($N_Q$)} & \textbf{Grid Step $\Delta\theta$ (rad)} & \textbf{Min. Amplitude $\sin^2(\Delta\theta)$} & \textbf{Estimated Gates} \\
\midrule
$3$ & $7$ & $0.392699$ & $0.146447$ & $12{,}068$ \\
$4$ & $15$ & $0.196350$ & $0.038060$ & $25{,}860$ \\
$\mathbf{5}$ & $\mathbf{31}$ & $\mathbf{0.098175}$ & $\mathbf{0.009607}$ & $\mathbf{53{,}444}$ \\
$6$ & $63$ & $0.049087$ & $0.002408$ & $108{,}612$ \\
$7$ & $127$ & $0.024544$ & $0.000602$ & $218{,}948$ \\
$8$ & $255$ & $0.012272$ & $0.000151$ & $439{,}620$ \\
$9$ & $511$ & $0.006136$ & $0.000038$ & $880{,}964$ \\
$10$ & $1{,}023$ & $0.003068$ & $0.000009$ & $1{,}763{,}652$ \\
\bottomrule
\end{tabular}
\end{table}

\subsection{Fault-Tolerant Threshold Implications}
On superconducting transmon architectures without quantum error correction, physical two-qubit gate fidelities are fundamentally bounded by decoherence and gate execution times:
\begin{equation}
D_{\mathrm{max}} \sim \frac{1}{\epsilon_{2q}} \approx \frac{1}{1.2 \times 10^{-2}} \approx 80 - 100 \text{ entangling layers}.
\end{equation}
As shown in Table \ref{tab:qae_scaling_registers}:
\begin{itemize}[leftmargin=*]
    \item At $n_E = 5$ (the baseline implemented in this work), the compiled circuit depth of $52{,}393$ exceeds the coherent threshold of unmitigated NISQ physical devices by more than two orders of magnitude. Consequently, full QPE-based QAE strictly requires fault-tolerant logical qubits supported by surface codes.
    \item To resolve ultra-rare industrial technosignatures ($p \sim 10^{-6}$, requiring $n_E \ge 11$ or iterative amplitude amplification techniques), the gate count exceeds $3 \times 10^6$ operations.
\end{itemize}
These empirical profiles justify the methodological division adopted in this thesis: validating ideal quadratic amplitude amplification on classical statevector emulators, while deploying compact, error-mitigated inference kernels with Zero-Noise Extrapolation for physical execution on intermediate-scale quantum processors.
